% Distribución íntegra por residencia del propietario cosmológico.
\section{Empalme geométrico entre contracción anisótropa, rebote torsional y región cosmológica descendiente}

El Big Bang puede entenderse como una distorsión de máxima intensidad entre estado,
rango y lector. No es una explosión ocurrida dentro de un espacio anterior. Es el
límite en el que la proyección que define extensión, distancia y duración empieza a
ser válida. Preguntar qué ocurrió «antes» utilizando el mismo reloj puede carecer de
sentido porque el reloj pertenece al régimen que está emergiendo. HMT conserva, sin
embargo, una genealogía anterior a esa publicación: ruta, frontera, orientación y
memoria.

La completación
\[
 1000=729+243+27+1
\]
muestra cómo una unidad compacta se distribuye en ambiente ternario, palabra
visible, cierre y unidad residual. No identifica por sí sola cuatro eras
cosmológicas. Proporciona una arquitectura de estratos capaz de distinguir el
estado profundo de su proyección. El mapa cosmológico toma la forma
\[
 \mathfrak R_{\rm cos}:
 \text{estado HMT de frontera}
 \dashrightarrow
 (M,g,e,\omega,\Psi).
\]
Su tarea es asignar rango, escala, conexión, contenido material y ley de evolución.
La tesis de la distorsión inicial afirma que la gran curvatura observada es la sombra
continua de una transición de ese mapa.

La torsión de Einstein--Cartan proporciona un mecanismo dinámico de rebote. Para una
materia con densidad de espín,
\[
 s^2\propto a^{-6}.
\]
La contribución torsional crece hacia volúmenes pequeños más deprisa que la materia
ordinaria. Una ecuación efectiva puede escribirse como
\[
 H^2
 =
 \frac{8\pi G}{3}
 \left(\rho_{\rm m}-\rho_{\rm spin}\right)
 -
 \frac{k}{a^2}
 +
 \frac{\Lambda}{3},
\]
con los factores de \(c\) restaurados según se empleen densidades de masa o de
energía. La igualdad
\[
 \rho_{\rm m}=\rho_{\rm spin}
\]
localiza el candidato de cancelación material--torsional. Hay rebote cuando, en una
escala no nula, la expresión completa del miembro derecho se anula,
\[
 H^2=0,
\]
y la ecuación de aceleración da \(\dot H>0\); los términos de curvatura y
\(\Lambda\) deben anularse, compensarse o incluirse expresamente en esa condición.
Bajo estas hipótesis, la contracción alcanza un mínimo y se invierte. El mecanismo
no introduce una presión arbitraria: procede de la ecuación de Cartan y de la
corriente de espín.

HMT interpreta ese umbral como cierre y reapertura de hoja. Durante la contracción,
la extensión visible se compacta y aumenta la memoria interna. En el punto de
retorno, la orientación publicada cambia y una nueva rama expansiva se hace
accesible. La singularidad deja de ser el único final posible porque la variable de
estado no se reduce al factor de escala. La ruta enriquecida puede continuar cuando
la carta cosmológica anterior pierde rango.

La misma geometría permite formular un universo hijo. Para un observador exterior,
una región en colapso queda tras un horizonte. Para la sección global, el cierre
puede abrir una frontera con orientación y tiempo propios. El empalme se describe
como un cobordismo o una transición entre componentes de hoja, no como la aparición
de materia desde la nada. La información que no se publica en el universo padre
puede convertirse en condición de frontera del universo hijo.

Esta imagen coordina tres fenómenos que suelen estudiarse por separado: colapso,
horizonte y origen cosmológico. La misma operación de cierre que forma una región
atrapada puede producir una reapertura tras alcanzar la escala torsional. La misma
memoria que protege la información exterior puede fijar las condiciones de la nueva
rama. Y el mismo cambio de orientación que define partícula y antipartícula a escala
microscópica puede adquirir una realización causal macroscópica.

La realización completa exige una solución global. Debe contener un exterior de
Einstein, una región de torsión no despreciable, un mínimo regular del factor de
escala, una continuación causal y un mapa de estados entre fronteras. Debe respetar
la conservación covariante y producir un observable: distribución de remanentes,
correlaciones no térmicas, modificación del espectro de ondas gravitatorias o una
huella primordial. Si la escala de rebote contradice nucleosíntesis, fondo cósmico o
límites de torsión, el mecanismo concreto queda descartado. La tesis de hoja no
exime a la dinámica; le proporciona el objeto que debe evolucionar.

\section{Regímenes elíptico, parabólico e hiperbólico del TRIT y su proyección sobre el sector oscuro}

El \(\TRIT\) no sólo orienta pasado, umbral y apertura. Sus tres valores determinan
las tres álgebras cuadráticas reales
\[
 \mathbb A_\tau
 =
 \mathbb R[X]/(X^2+\bar\tau).
\]
Los generadores satisfacen
\[
 \bar\tau=+1:
 \quad J^2=-I,
 \qquad
 \bar\tau=0:
 \quad J^2=0,
 \qquad
 \bar\tau=-1:
 \quad J^2=+I.
\]
Son los regímenes elíptico, parabólico e hiperbólico. La geometría modelo
\[
 g_{s,\rho}
 =
 \dd r^2+S_{s,\rho}(r)^2\dd\theta^2
\]
se obtiene con
\[
 S_{+,\rho}(r)=\rho^{-1}\sin(\rho r),
 \qquad
 S_{0,\rho}(r)=r,
 \qquad
 S_{-,\rho}(r)=\rho^{-1}\sinh(\rho r),
\]
y posee
\[
 K_G=s\rho^2.
\]
Las tres curvaturas no se atribuyen simultáneamente a la misma métrica en un mismo
punto. Son hojas coordinadas de un estado cuya carta física selecciona el régimen
realizado.

El espacio de de Sitter demuestra de forma exacta cómo una sola geometría
cuadridimensional puede admitir tres presentaciones espaciales:
\[
 \dd s^2
 =
 -\dd t^2+a_s(t)^2\dd\Sigma_s^2,
 \qquad s\in\{+1,0,-1\}.
\]
Las funciones
\[
 \begin{aligned}
 a_+(t)&=H^{-1}\cosh(Ht),
 &
 a_0(t)&=\eexp^{Ht},\\
 a_-(t)&=H^{-1}\sinh(Ht),
 &
 t&>0,
 \end{aligned}
\]
satisfacen
\[
 \frac{\ddot a_s}{a_s}=H^2,
\qquad
 \left(\frac{\dot a_s}{a_s}\right)^2
 +
 \frac{s}{a_s^2}
 =
 H^2,
\qquad
 R=12H^2.
\]
Una foliación cerrada, una plana y una abierta describen el mismo espacio-tiempo.
Planitud espacial local, sección hiperbólica y aceleración positiva no son signos
incompatibles. Pueden ser distintas hojas de una misma geometría global.

Este hecho proporciona un modelo continuo para la intuición HMT. El valor publicado
en una hoja no contiene necesariamente toda la curvatura del estado. Un observador
que trabaja en la foliación plana puede reconstruir como fuente adicional el efecto
que procede de incidencia de frontera o de otra presentación. La materia oscura y
la energía oscura reciben así una interpretación unificada, pero diferenciada.

La materia oscura aparente puede corresponder a respuesta gravitatoria de estados
que no poseen representante puntual en la hoja visible. Su masa se manifiesta en
órbitas y lente porque la geometría completa transporta su incidencia, aunque el
lector electromagnético no publique una partícula ordinaria. La energía oscura
aparente puede corresponder a la respuesta global de apertura y curvatura de hoja,
publicada como aceleración del factor de escala. Una pertenece principalmente a la
distribución de retención no visible; la otra, a la evolución de la frontera
cosmológica.

La bisagra
\[
 \frac{G_{\rm tor}}{I_{\rm tor}}=1,
 \qquad
 \frac{G_{\rm av}}{I_{\rm av}}
 =
 \frac{\piHMT}{\phiHMT^2}
\]
indica dónde buscar ese residuo. La razón marcada es una relación estructural entre
lecturas, no un porcentaje cosmológico que pueda identificarse directamente con una
densidad observada. Para convertirse en teoría del sector oscuro debe integrarse
sobre la geometría y producir un tensor efectivo
\(T^{\rm eff}_{\mu\nu}\), sus ecuaciones de continuidad y una función de
transferencia.

La exigencia es especialmente fuerte porque una sola regla debe atravesar varias
escalas. En galaxias ha de relacionar densidad bariónica, velocidad circular y lente
sin introducir un parámetro nuevo para cada objeto. En cúmulos debe coordinar gas,
lente débil y dinámica. En cosmología debe predecir \(H(z)\), distancias,
crecimiento de estructura y potenciales de lente. Si la misma geometría multinoja
produce esos perfiles con parámetros fijados antes de consultar cada sistema,
HMT--MD habrá sustituido dos sectores desconocidos por una sola arquitectura. Si
requiere ajustes independientes, la interpretación no habrá alcanzado su objetivo.

La triple foliación impone además una disciplina conceptual. No toda curvatura
hiperbólica es una señal de energía oscura, ni toda planitud local elimina la
curvatura global. La carta, el dominio y la escala deben declararse. El parámetro
\(H\) de de Sitter tampoco queda seleccionado sólo por la existencia de las tres
foliaciones. La matemática prueba su compatibilidad; el transductor cosmológico debe
derivar la escala física y su evolución.

\section{Orientación temporal del fondo cosmológico y conservación de memoria bajo expansión}

El cierre dodecafásico produce una plantilla angular muy restringida. Sea
\[
 \phi_k=\phi_0+\frac{2\pi k}{12},
 \qquad k=0,\ldots,11,
\]
y
\[
 w_k
 =
 \frac1{12}
 \left[
 1+2\epsilon\cos(3\phi_k-\delta)
 \right],
 \qquad
 |\epsilon|\leq\frac12.
\]
Los pesos son no negativos y suman uno. Su transformada discreta tiene soporte sólo
en
\[
  m=0,\pm3\quad(\operatorname{mod}\ 12).
\]
Al proyectar sobre multipolos bajos, para \(\ell=2\) sólo puede sobrevivir
\(m=0\); para \(\ell=3\), \(m=0,\pm3\). La familia \(C_{12}\) convierte la rueda
dodecafásica en una envolvente angular exacta.

La importancia física reside en su restrictividad. El eje, la fase \(\delta\), la
amplitud \(\epsilon\), la máscara celeste, el estimador y el conjunto de multipolos
deben fijarse antes de examinar el mapa que se utilizará como contraste. Después se
comparan datos de frecuencias y misiones independientes con simulaciones isotrópicas
y con modelos de contaminación. Una señal estable en el soporte previsto reforzaría
la genealogía dodecafásica. Una selección posterior del eje o una señal dominante
fuera de \(m=0,\pm3\) excluirían esa realización concreta.

La flecha cosmológica del tiempo utiliza el mismo lector ternario que organiza las
hojas:
\[
 +1
 \longmapsto
 \text{retorno, memoria, pasado},
\]
\[
 0
 \longmapsto
 \text{umbral, presente},
\]
\[
 -1
 \longmapsto
 \text{apertura bidireccional, futuro}.
\]
El presente no es un punto que se desplaza sobre una recta previamente trazada. Es
el umbral en el que una apertura admisible se actualiza y empieza a conservarse como
memoria. El pasado es el sector de retornos ya inscritos; el futuro es el conjunto
de continuaciones compatibles que todavía no han cerrado.

La evolución completa puede ser reversible mientras la historia publicada posee
orientación. La asimetría aparece en la preparación, en la selección de frontera y
en el lector que abandona parte de la fibra. La entropía no crea el tiempo; mide la
información que una presentación deja de distinguir. Para una acción de historias
\(I(\gamma)\), la medida
\[
 \mathbb P_\beta(\gamma)
 =
 \frac{\eexp^{-\beta I(\gamma)}}{Z_\beta}
\]
minimiza
\[
 \mathbb E_{\mathbb P}[I]
 -
 \beta^{-1}H(\mathbb P).
\]
En el límite de baja temperatura domina la interacción mínima; a temperatura
finita interviene la multiplicidad de rutas. La acción y la entropía son dos
componentes de una misma selección, no dos causas rivales.

Una historia de ida y retorno puede purificarse como
\[
 |\Omega_\beta\rangle
 =
 \sum_\gamma
 \sqrt{\mathbb P_\beta(\gamma)}
 |\gamma\rangle_{\rm ida}
 |\mathcal C\gamma\rangle_{\rm vuelta}.
\]
La entropía reducida es \(H(\mathbb P_\beta)\). El universo observado puede poseer
una flecha porque accedemos a una rama, mientras el estado ampliado conserva la
conjugación. Esta imagen enlaza la orientación microscópica de partícula y
antipartícula con la asimetría macroscópica sin afirmar que ambas escalas sean
idénticas.

La cosmología HMT--MD queda así organizada por una misma operación fundamental:
plegar, conservar y desplegar. La materia es una sección que retiene espacio; el
horizonte pliega una presentación en otra hoja; el rebote despliega una rama después
de la compactación; el sector oscuro es la respuesta de información geométrica que
la hoja visible no publica; y el tiempo es la orientación de apertura y retorno.
Cada uno de estos enunciados adquiere fuerza física al compartir el mismo dominio de
estado, no por formar una colección de analogías.

La teoría se compromete con consecuencias claras. El exterior del agujero negro debe
seguir la relatividad en su régimen probado y la inversión debe aportar una
continuación y una señal. El rebote debe fijar su escala mediante la densidad de
espín y respetar nucleosíntesis y fondo cósmico. La función de transferencia
multinoja debe explicar conjuntamente dinámica, lente y expansión. La plantilla
\(C_{12}\) debe sobrevivir a un contraste prospectivo. La flecha temporal debe
respetar reversibilidad global y coste de borrado. Estas condiciones no rebajan la
tesis cosmológica: son el camino por el que una arquitectura matemática se convierte
en una física del universo.

